\documentclass[11pt,a4paper]{article}

\usepackage[utf8]{inputenc}
\usepackage[T1]{fontenc}
\usepackage{textcomp}
\usepackage[french]{babel}
\usepackage{lmodern}
\usepackage[a4paper,left=3cm,right=3cm,top=2.8cm,bottom=2.8cm,headheight=14pt]{geometry}
\usepackage{booktabs}
\usepackage{tabularx}
\usepackage{array}
\usepackage{enumitem}
\usepackage{fancyhdr}
\usepackage{titlesec}
\usepackage[hidelinks]{hyperref}

\hypersetup{pdftitle={Recommandation de chaussures par photo},
            pdfauthor={}, pdfsubject={}, pdfkeywords={}, pdfcreator={}, pdfproducer={}}

% --- Titres : sobres, sans couleur ---
\titleformat{\section}{\large\bfseries}{\thesection.}{0.6em}{}
\titlespacing*{\section}{0pt}{1.8em}{0.7em}
\titleformat{\subsection}{\normalsize\bfseries}{\thesubsection}{0.6em}{}
\titlespacing*{\subsection}{0pt}{1.3em}{0.4em}

% --- En-tête / pied ---
\pagestyle{fancy}
\fancyhf{}
\renewcommand{\headrulewidth}{0.4pt}
\fancyhead[L]{\footnotesize Recommandation de chaussures par photo}
\fancyhead[R]{\footnotesize\thepage}

\renewcommand{\arraystretch}{1.3}
\newcolumntype{L}[1]{>{\raggedright\arraybackslash}p{#1}}

% --- Bloc de remarque : filets simples, pas d'encadré coloré ---
\newenvironment{remarque}
  {\par\vspace{0.9em}\noindent\rule{\linewidth}{0.4pt}\par\vspace{0.5em}
   \begingroup\leftskip=1em\rightskip=1em}
  {\par\endgroup\vspace{0.5em}\noindent\rule{\linewidth}{0.4pt}\par\vspace{0.9em}}

\setlist[itemize]{leftmargin=1.3em,itemsep=0.3em,topsep=0.4em}
\setlist[enumerate]{leftmargin=1.5em,itemsep=0.4em,topsep=0.4em}

\begin{document}

% =====================================================================
\thispagestyle{empty}
\vspace*{2cm}

\noindent
{\footnotesize NOTE TECHNIQUE ET BUDGÉTAIRE}

\vspace{0.6cm}
\noindent
{\LARGE\bfseries Moteur de recommandation de chaussures par photo}

\vspace{0.5cm}
\noindent
{\large Fonctionnement, technologies employées et coûts d'hébergement}

\vspace{0.4cm}
\noindent\rule{\linewidth}{0.4pt}

\vspace{1.2cm}
\noindent
Le client de la boutique en ligne envoie une photo. Le système détermine s'il
s'agit d'une chaussure ou d'une tenue, puis propose des articles du catalogue,
en stock, avec lien vers la fiche produit et ajout au panier. Aucune image
n'est générée : chaque proposition est l'une de vos propres photos produit.

\vspace{0.8cm}
\noindent
Le présent document décrit le fonctionnement du système, les briques
techniques qui le composent, la configuration serveur nécessaire, et le coût
de fonctionnement mensuel selon que l'on conserve ou non le service d'IA de
langage.

\vspace{2cm}
\noindent
{\footnotesize 31 juillet 2026}

\newpage

{\small\tableofcontents}

\newpage

% =====================================================================
\section{Les deux modes de fonctionnement}

Le système répond à une question simple posée par le visiteur : \emph{j'aime
cette photo, qu'avez-vous qui y ressemble, ou qui irait avec ?} Deux modes
existent, choisis automatiquement selon la photo reçue.

\vspace{0.5em}
\begin{tabularx}{\textwidth}{L{3.2cm} X}
\toprule
\textbf{Photo reçue} & \textbf{Ce que voit le client} \\
\midrule
Une chaussure &
Une grille de chaussures visuellement proches, classées par ressemblance.
Chaque résultat est un produit réel du catalogue, en stock, cliquable. \\
Une tenue &
Trois chaussures qui complètent la tenue, chacune accompagnée d'une phrase de
conseil rédigée par une IA de langage (\emph{« ce derby noir structure ce
costume marine pour le bureau »}). \\
\bottomrule
\end{tabularx}

\vspace{1em}
Ces deux modes n'ont pas le même coût de fonctionnement, et c'est le point
déterminant pour le budget.

Le mode chaussure vers chaussures similaires est \textbf{entièrement gratuit à
l'usage}. Tout le calcul se fait sur votre serveur, sans appel à un service
extérieur facturé. Que vos clients envoient cent ou cent mille photos par mois,
la facture ne change pas.

Le mode tenue vers chaussures assorties est \textbf{le seul poste facturé à la
photo}, car c'est lui qui fait appel à l'IA de langage pour choisir les trois
chaussures et rédiger les phrases de conseil. Ce coût reste faible : environ
3~€ pour mille photos de tenues traitées.

% =====================================================================
\section{Comment fonctionne la partie gratuite}
\label{sec:sansllm}

C'est la partie qui fait la majorité du travail, et c'est celle qui ne coûte
rien à l'usage. Elle mérite d'être comprise, parce qu'elle explique pourquoi le
budget de ce projet est aussi contenu.

\subsection{Le principe : une empreinte visuelle par photo}

Imaginons que l'on puisse résumer n'importe quelle photo par une longue suite de
chiffres, une sorte d'empreinte visuelle comparable à une empreinte digitale.
Cette empreinte ne décrit pas la photo avec des mots ; elle capture sa signature
visuelle : les formes, les matières, les couleurs, l'allure d'ensemble.

La propriété utile est la suivante : deux photos qui se ressemblent produisent
deux empreintes proches l'une de l'autre. Une botte noire en cuir photographiée
de face et la même botte photographiée de trois quarts donnent des empreintes
très voisines. Une basket blanche donnera une empreinte nettement plus éloignée.

Un modèle de vision par ordinateur, installé une fois pour toutes sur le
serveur, calcule ces empreintes. Il ne comprend pas au sens humain : il compare.
Et il fonctionne en local, sans connexion à un service payant.

\subsection{Ce qui se passe concrètement}

\begin{enumerate}
  \item \textbf{Une fois, à la mise en service.} On parcourt le catalogue.
  Pour chaque photo de chaque produit, on calcule l'empreinte et on la range
  dans une base spécialisée. Un même modèle photographié sous trois angles donne
  trois empreintes qui pointent toutes vers la même fiche produit : c'est ce qui
  rend la recherche insensible à l'angle de prise de vue.

  \item \textbf{À chaque photo envoyée.} On calcule l'empreinte de la photo du
  client, exactement de la même façon, puis on cherche les empreintes les plus
  proches dans la base. Environ un dixième de seconde.

  \item \textbf{On filtre et on affiche.} On ne retient que les produits en
  stock, on évite de présenter six fois le même modèle dans six coloris, et on
  affiche une grille cliquable.
\end{enumerate}

\subsection{Le même mécanisme rend trois autres services}

\begin{itemize}
  \item \textbf{Reconnaître le type de photo.} Le système compare l'empreinte
  reçue à celle du concept « une chaussure » et à celle du concept « une tenue
  portée ». La plus proche l'emporte. Si les deux sont trop serrées, la photo
  est jugée ambiguë et la boutique affiche les deux résultats sous forme
  d'onglets.

  \item \textbf{Décrire les produits.} Le même mécanisme attribue à chaque
  chaussure son type (botte, derby, basket\ldots), sa matière, ses couleurs
  dominantes et un niveau d'habillement de 1 à 5. Ces étiquettes sont calculées
  une fois à l'indexation, puis stockées.

  \item \textbf{Préparer le travail de l'IA.} En mode tenue, ce sont ces filtres
  gratuits qui ramènent le catalogue entier à une trentaine de candidats
  pertinents. L'IA payante ne voit jamais tout le catalogue : elle reçoit une
  courte liste déjà triée. C'est ce qui maintient son coût à un niveau
  négligeable.
\end{itemize}

% =====================================================================
\section{Ce qu'implique le retrait de l'IA de langage}
\label{sec:retrait}

Il est important d'être précis sur ce point, car la réponse est moins favorable
qu'on pourrait le supposer.

Dans l'état actuel du logiciel, l'IA de langage ne se contente pas de rédiger
les phrases de conseil : c'est elle qui effectue la sélection finale des trois
chaussures. Si la clé d'accès au service est retirée, l'appel échoue et l'erreur
remonte jusqu'à la boutique.

\begin{remarque}
\noindent
\textbf{Comportement réel sans clé d'accès}, vérifié dans le code :

\vspace{0.4em}
\begin{tabularx}{\linewidth}{L{4.2cm} X}
Photo d'une chaussure & Fonctionne normalement. Aucune dégradation. \\
Photo d'une tenue & \textbf{Message d'erreur} affiché au client. \\
Photo ambiguë & \textbf{Message d'erreur}, alors même que les chaussures
similaires avaient été calculées correctement. \\
\end{tabularx}

\vspace{0.6em}
\noindent
Autrement dit, retirer purement et simplement l'IA ne dégrade pas seulement le
mode tenue : cela produit une erreur visible pour le client dès qu'une photo
n'est pas clairement identifiée comme une chaussure.
\end{remarque}

Si l'objectif est de supprimer entièrement le coût variable, il ne faut donc pas
se contenter de retirer la clé : il faut prévoir un développement
complémentaire, à savoir intercepter l'échec de l'appel et faire retomber le
mode tenue sur un classement purement visuel. Le client verrait alors des
chaussures proposées, sans phrase de conseil, et aucune erreur ne s'afficherait.

Ce développement représente environ une journée de travail. Il n'est pas
réalisé aujourd'hui, et son coût doit être mis en regard de l'économie
attendue, de l'ordre de quelques euros par mois (voir §\ref{sec:comparatif}).

% =====================================================================
\newpage
\section{Les briques technologiques}

Le système assemble six composants. Cinq sont des logiciels libres, gratuits et
sans redevance, installés sur votre serveur. Un seul est un service extérieur
facturé à l'usage.

\vspace{0.6em}
\begin{tabularx}{\textwidth}{L{2.9cm} X L{2.6cm}}
\toprule
\textbf{Composant} & \textbf{Rôle} & \textbf{Licence} \\
\midrule

Modèle de vision &
Calcule les empreintes visuelles décrites au §\ref{sec:sansllm}. Assure la
reconnaissance du type de photo, la description des produits et la mesure de
ressemblance. Chargé une fois au démarrage et gardé en mémoire. &
Gratuit \\

Base d'empreintes &
Stocke une empreinte par photo produit et retrouve les plus proches en quelques
millisecondes. &
Gratuit \\

Base de données &
Conserve les informations lisibles de chaque produit : titre, prix, lien,
disponibilité, étiquettes calculées. Permet de filtrer par style, saison et
prix. &
Gratuit \\

Serveur applicatif &
Reçoit les photos, orchestre les étapes, renvoie les résultats. Comprend les
garde-fous : limite de 5~Mo par image, rejet des fichiers non-image, limitation
du nombre de requêtes par visiteur. &
Gratuit \\

Module boutique &
Le bloc visible par vos clients : bouton d'envoi, aperçu, grille de résultats,
onglets en cas de photo ambiguë, ajout au panier. S'installe dans le thème via
l'éditeur de la plateforme. &
Gratuit \\

\midrule
IA de langage &
Le seul poste facturé à l'usage. Intervient uniquement en mode tenue : reçoit la
trentaine de candidats pré-sélectionnés, choisit les trois meilleurs et rédige
une phrase pour chacun. Un appel par photo de tenue. &
À l'usage \\
\bottomrule
\end{tabularx}

\vspace{1em}
Deux conséquences méritent d'être signalées.

Le fournisseur d'IA est interchangeable. Passer à un autre modèle, moins cher ou
plus performant, est un changement d'une ligne de configuration, sans
réécriture.

Aucune carte graphique n'est nécessaire. Le modèle de vision a été choisi pour
tourner sur un processeur ordinaire, ce qui permet d'héberger l'ensemble sur un
serveur à moins de dix euros par mois plutôt que sur une machine spécialisée à
plusieurs centaines d'euros.

% =====================================================================
\section{Configuration d'hébergement}

\subsection{Ce n'est pas la taille du catalogue qui dimensionne}

Point contre-intuitif : pour un catalogue de moins de 500 produits, les données
occupent une place négligeable, quelques mégaoctets d'empreintes et de fiches.
Ce n'est donc pas le catalogue qui détermine le serveur à louer.

Ce qui dimensionne, c'est la mémoire vive. Le modèle de vision doit rester
chargé en permanence pour que les réponses soient instantanées, ce qui mobilise
de l'ordre de 1~Go, auxquels s'ajoutent les deux bases de données et le système.
Il faut également de l'espace disque pour l'image logicielle, qui embarque la
bibliothèque de calcul et les poids du modèle.

\vspace{0.6em}
\begin{tabularx}{\textwidth}{L{3.8cm} L{3.2cm} X}
\toprule
\textbf{Ressource} & \textbf{Recommandation} & \textbf{Justification} \\
\midrule
Mémoire vive & 8 Go & 4 Go suffisent en théorie mais laissent peu de marge ;
8 Go écarte tout risque d'arrêt en période de pointe. \\
Processeur & 2 à 4 cœurs & Le calcul d'empreinte prend environ un dixième de
seconde par image. \\
Disque & 40 à 80 Go & Essentiellement pour l'image logicielle. \\
Carte graphique & Aucune & Volontairement écarté à la conception. \\
Système & Docker & L'ensemble se lance par une seule commande. \\
\bottomrule
\end{tabularx}

\vspace{1em}
Les offres d'hébergement gratuites ou d'entrée de gamme (256~Mo à 1~Go de
mémoire) sont insuffisantes : le service ne démarrera pas. Il faut un serveur
privé virtuel classique, ce qui reste peu coûteux.

\subsection{Deux formules}

\begin{tabularx}{\textwidth}{L{3.6cm} X L{2.4cm}}
\toprule
\textbf{Formule} & \textbf{Description} & \textbf{Par mois} \\
\midrule
Serveur privé virtuel &
Une machine louée chez un hébergeur européen, sur laquelle tout tourne.
Configuration 4 cœurs / 8 Go / 80 Go. Nécessite une mise en place initiale et
des mises à jour occasionnelles. \emph{Formule retenue pour les calculs qui
suivent.} &
8 à 11 € \\
\addlinespace
Plateforme managée &
L'hébergeur assure les mises à jour, les sauvegardes et la surveillance. Plus
confortable, sensiblement plus cher, souvent facturé par composant. &
25 à 60 € \\
\bottomrule
\end{tabularx}

\vspace{0.8em}
\noindent
{\small Tarifs relevés le 31 juillet 2026 auprès d'hébergeurs européens pour une
configuration 4 cœurs / 8 Go / 80 Go, soit 8,49~€ à 10,49~€ hors taxes par mois
selon le type de processeur. Ces tarifs ont augmenté en juin 2026 et sont à
revérifier au moment de la commande. Les chiffres des sections suivantes
retiennent une hypothèse de \textbf{9~€ par mois}, à laquelle s'ajoute environ
1~€ pour le nom de domaine et le certificat de sécurité.}

% =====================================================================
\newpage
\section{Coût de l'IA de langage}
\label{sec:llm}

\subsection{Mode de calcul}

L'IA est facturée à la quantité de texte échangée, et non au nombre de
requêtes. Elle ne reçoit ni la photo, ni le catalogue complet : elle reçoit une
courte liste de texte, une trentaine de chaussures déjà pré-sélectionnées par la
partie gratuite, réduites à leurs seules caractéristiques utiles (type, matière,
couleurs, prix). Elle renvoie trois identifiants et trois phrases. L'échange est
donc volontairement minuscule.

\vspace{0.6em}
\begin{tabularx}{\textwidth}{X L{3.4cm} L{2.4cm}}
\toprule
\textbf{Élément} & \textbf{Volume} & \textbf{Coût} \\
\midrule
Texte envoyé (tenue + 30 candidats) & environ 2\,100 unités & 0,0019 € \\
Texte renvoyé (3 choix + phrases) & environ 200 unités & 0,0009 € \\
\midrule
\textbf{Total par photo de tenue} & & \textbf{0,003 €} \\
\bottomrule
\end{tabularx}

\vspace{1em}
Soit environ 3~€ pour mille photos de tenues traitées. Les photos de chaussures
ne déclenchent aucun appel et ne coûtent rien.

\subsection{Projection sur trois niveaux d'activité}

L'hypothèse retenue est que 40~\% des photos envoyées sont des tenues et 60~\%
des chaussures, répartition prudente pour une boutique de chaussures où l'usage
« je photographie une chaussure qui me plaît » domine généralement. Le serveur
est compté à 9~€ par mois.

\vspace{0.6em}
\begin{tabularx}{\textwidth}{L{3.4cm} L{2.6cm} L{2.8cm} X}
\toprule
\textbf{Photos / mois} & \textbf{Dont tenues} & \textbf{Coût IA} &
\textbf{Serveur + IA} \\
\midrule
500 & 200 & 0,60 € & 9,60 € \\
2\,000 & 800 & 2,40 € & 11,40 € \\
10\,000 & 4\,000 & 12,00 € & 21,00 € \\
\bottomrule
\end{tabularx}

\vspace{1em}
Le coût de l'IA reste marginal même à fort trafic : multiplier l'activité par
vingt ne multiplie la facture totale que par deux, parce que le serveur, poste
principal, est un forfait fixe.

Ce projet ne présente donc pas de risque de dérapage budgétaire lié au succès.
C'est une conséquence directe du choix d'architecture : faire porter l'essentiel
du travail par des logiciels gratuits tournant en local, et réserver l'IA payante
à une seule étape, sur une liste déjà réduite.

% =====================================================================
\section{Comparatif : avec et sans IA de langage}
\label{sec:comparatif}

\begin{tabularx}{\textwidth}{L{4.4cm} X X}
\toprule
& \textbf{Avec IA} & \textbf{Sans IA} \\
\midrule
Photo d'une chaussure & Chaussures similaires & Identique \\
Photo d'une tenue & 3 chaussures + conseil & Erreur, sauf développement
complémentaire \\
Photo ambiguë & Deux onglets & Erreur, sauf développement complémentaire \\
Reconnaissance du type & Incluse, gratuite & Incluse, gratuite \\
Description du catalogue & Incluse, gratuite & Incluse, gratuite \\
\midrule
Serveur et domaine & 10 € / mois & 10 € / mois \\
IA, selon le trafic & 0,60 à 12 € / mois & 0 € \\
\midrule
\textbf{Total mensuel} & \textbf{10,60 à 22 € / mois} & \textbf{10 € / mois} \\
\bottomrule
\end{tabularx}

\vspace{1.2em}
L'écart entre les deux formules correspond exactement au poste IA, soit
\textbf{0,60 à 12~€ par mois} selon le trafic.

Cette économie est à mettre en regard de deux éléments. D'une part, le mode
tenue est précisément l'argument qui distingue ce système d'un moteur de
recherche classique. D'autre part, le retrait de l'IA suppose une journée de
développement complémentaire (§\ref{sec:retrait}) pour éviter d'afficher des
erreurs aux clients, ce qui annule à lui seul plusieurs années d'économie au
niveau de trafic le plus bas.

\vspace{0.5em}
\noindent
\textbf{Recommandation : conserver l'IA de langage.} Si le budget doit être
contenu, les leviers pertinents sont ailleurs :

\begin{itemize}
  \item conserver la description gratuite du catalogue plutôt que la variante
  payante (§\ref{sec:ponctuels}) ;
  \item conserver le modèle d'IA économique configuré par défaut, qui suffit
  largement pour rédiger une phrase de conseil ;
  \item choisir la formule serveur privé virtuel plutôt que la plateforme
  managée, si vous disposez d'un interlocuteur technique.
\end{itemize}

% =====================================================================
\newpage
\section{Coûts ponctuels de mise en service}
\label{sec:ponctuels}

\begin{tabularx}{\textwidth}{L{4.6cm} X L{2.2cm}}
\toprule
\textbf{Poste} & \textbf{Détail} & \textbf{Coût} \\
\midrule
Indexation initiale &
Parcours de tous les produits pour calculer et stocker les empreintes. Pour
moins de 500 produits : moins de cinq minutes de traitement. &
0 € \\
\addlinespace
Description enrichie \newline (option) &
Variante payante faisant décrire chaque produit par une IA plutôt que par le
modèle de vision gratuit. Étiquettes légèrement plus fines, à refacturer à
chaque réindexation complète. Non recommandé. &
0,50 € \\
\addlinespace
Mise en ligne &
Location du serveur, installation, mise en place du bloc dans le thème,
connexion au catalogue. &
Prestation \\
\bottomrule
\end{tabularx}

\vspace{1em}
Le catalogue se met ensuite à jour automatiquement : toute modification d'un
produit dans l'administration de la boutique (prix, rupture de stock, nouvelle
photo) est répercutée dans l'index dans la foulée, sans intervention et sans
coût.

% =====================================================================
\section{Synthèse}

\begin{remarque}
\noindent
\textbf{Budget de fonctionnement, catalogue de moins de 500 produits}

\vspace{0.5em}
\begin{tabularx}{\linewidth}{X L{4.4cm}}
Serveur privé virtuel & 9 € / mois \\
Nom de domaine et certificat & 1 € / mois \\
IA de langage, selon le trafic & 0,60 à 12 € / mois \\
Licences logicielles & 0 € \\
\midrule
\textbf{Total} & \textbf{10,60 à 22 € / mois} \\
\end{tabularx}
\end{remarque}

Trois points à retenir.

\begin{enumerate}
  \item \textbf{Cinq des six composants sont libres et gratuits.} Aucune
  redevance logicielle, aucun abonnement par utilisateur, aucun coût qui
  augmente avec la taille du catalogue.

  \item \textbf{La recherche par ressemblance visuelle ne coûte rien à
  l'usage}, parce qu'elle repose sur un calcul local et non sur un service
  extérieur. C'est le mode le plus sollicité sur une boutique de chaussures, et
  c'est celui qui est gratuit.

  \item \textbf{Le seul poste variable est plafonné par construction.} L'IA
  n'intervient qu'une fois par photo de tenue, sur une liste de trente candidats
  déjà filtrés. Même à dix mille photos par mois, elle reste du même ordre de
  grandeur que le coût du serveur.
\end{enumerate}

% =====================================================================
\vspace{1.5em}
\section*{Hypothèses et sources}
\addcontentsline{toc}{section}{Hypothèses et sources}

{\small
\begin{itemize}
  \item Tarifs d'hébergement relevés le 31 juillet 2026 auprès d'hébergeurs
  européens, configuration 4 cœurs / 8 Go de mémoire / 80 Go de disque, hors
  taxes, soit 8,49~€ à 10,49~€ par mois. Ces tarifs ont fait l'objet d'une
  hausse en juin 2026 et sont à revérifier au moment de la commande. Les
  projections retiennent 9~€ par mois.

  \item Tarif de l'IA de langage : 1~\$ par million d'unités reçues et 5~\$ par
  million d'unités produites, tarif public du modèle économique configuré par
  défaut, relevé le 31 juillet 2026. Conversion appliquée : 1~\$ = 0,92~€.

  \item Volumes d'échange estimés à partir du contenu réel des requêtes du
  système, soit environ 2\,100 unités reçues et 200 unités produites par photo
  de tenue.

  \item Répartition supposée du trafic : 40~\% de photos de tenues, 60~\% de
  photos de chaussures. Une proportion de tenues plus faible réduirait
  proportionnellement le coût variable.

  \item Comportement décrit au §\ref{sec:retrait} en l'absence de clé d'accès :
  vérifié par lecture du code source, et non estimé.

  \item Catalogue de référence : moins de 500 produits, plusieurs photos par
  produit. Un catalogue plus important allongerait l'indexation initiale mais ne
  modifierait ni la configuration serveur ni le coût par photo.
\end{itemize}
}

\end{document}
